\documentclass[11pt,reqno]{amsart}
\usepackage[T1]{fontenc}
\usepackage{lmodern}
\usepackage{microtype}
\usepackage{amsmath,amssymb,mathtools}
\usepackage{booktabs}
\usepackage[colorlinks=true,linkcolor=blue,citecolor=blue,urlcolor=blue]{hyperref}

\newtheorem{theorem}{Theorem}[section]
\newtheorem{lemma}[theorem]{Lemma}
\newtheorem{proposition}[theorem]{Proposition}
\newtheorem{corollary}[theorem]{Corollary}
\newtheorem{conjecture}[theorem]{Conjecture}
\newtheorem{hypothesis}[theorem]{Hypothesis}
\newtheorem*{principle}{Principle}
\newtheorem{introthm}{Theorem}
\renewcommand{\theintrothm}{\Alph{introthm}}
\theoremstyle{definition}
\newtheorem{definition}[theorem]{Definition}
\newtheorem{remark}[theorem]{Remark}
\newtheorem{evidence}[theorem]{Evidence}

\newcommand{\Z}{\mathbb Z}
\newcommand{\Q}{\mathbb Q}
\newcommand{\Pp}{\mathsf P}
\newcommand{\cP}{\mathcal P}
\newcommand{\cV}{\mathcal V}
\newcommand{\cF}{\mathfrak F}
\newcommand{\Se}{S_{\mathrm e}}
\newcommand{\Sa}{S_{\mathrm a}}
\newcommand{\dist}{\operatorname{dist}}
\newcommand{\Fib}{\operatorname{Fib}}
\newcommand{\lc}{\operatorname{lc}}
\newcommand{\leg}[2]{\left(\frac{#1}{#2}\right)}
\newcommand{\status}[1]{\textsc{#1}}

% TODO(submission): author metadata (parent to settle).
\title[The pointwise signed-graph approach to Erd\H{o}s--Straus is blocked]{Formally generic primes block the pointwise signed-graph approach to the Erd\H{o}s--Straus conjecture}
\author{Anonymous}
\date{Draft, 1 October 2026}

\begin{document}

\begin{abstract}
For a prime $p=4t+1$ the signed integer solutions of $4/p=1/x+1/y+1/z$
form a finite graph, two solutions being adjacent when they share a
denominator.  The \emph{seed} $(t,-2pt,-2pt)$ is available for every such
$p$, and the \emph{seed-component conjecture} asserts that its component
always contains an all-positive solution; it would imply the
Erd\H{o}s--Straus conjecture for $p\equiv1\pmod4$.  We explain why this
pointwise programme cannot succeed by any argument that uses only
congruence data and the shape of factorisations.  Unconditionally, we
classify all escapes from the seed of length at most three, show that every
prime outside Mordell's six classes modulo $840$ escapes in two moves, and
bound the exceptional sets by $N(\log N)^{-11/2}$ (distance $>2$) and
$N(\log N)^{-10}$ (distance $>5$).  Under Schinzel's Hypothesis~H the seed
distance is unbounded.  Moreover, the whole seed component is free of positive
vertices for infinitely many primes.  This was proved first, and in a
stronger form, by a companion computation assuming Dickson's conjecture for
$159$ linear forms only; that construction also exhibits a positive
solution outside the sterile component.  We give an independent
confirmation from Hypothesis~H for $6402$ explicit polynomials of degree at
most two, via a certified finite closure of $7883$ formal vertices.  A
common mechanism lies behind these results, the Elsholtz--Tao
``odd-square'' remark and Schinzel's polynomial-identity theorem: at a
formally generic prime every number the argument meets factors formally, and
by the Bright--Loughran character dichotomy no formal solution is positive.
We also record further obstructions: sterile components larger than the
seed component, the triviality of large $p$-free buckets, and parity
obstructions.  The Erd\H{o}s--Straus conjecture itself is untouched.
\end{abstract}

\maketitle



%======================================================================
\section{Introduction}\label{sec:intro}

The Erd\H{o}s--Straus conjecture (ES) asserts that for every integer $n\ge2$
\begin{equation}\label{eq:ES}
 \frac4n=\frac1x+\frac1y+\frac1z
\end{equation}
has a solution in positive integers.  It suffices to treat primes, and
Mordell's identities settle all primes outside the six classes
$p\equiv1,121,169,289,361,529\pmod{840}$ \cite{Mordell}.  The conjecture has
been verified numerically to $10^{18}$ \cite{MD}.  It is open; no published
argument proves it.

\subsection{The signed graph and the seed}
Allowing \emph{signed} denominators changes the problem completely: signed
solutions exist for every odd $n$ (Jaroma's identity, see \cite[\S1]{BL}),
and for fixed $n$ there are only finitely many of them \cite[Lemma~3.10]{BL}.
Throughout, $p=4t+1$ is a prime.

\begin{definition}\label{def:graph}
A \emph{vertex} is a multiset $\{x,y,z\}$ of nonzero integers with
$4/p=1/x+1/y+1/z$.  It is \emph{positive} if $x,y,z>0$ and
\emph{nonpositive} otherwise.  The \emph{signed graph} $G_p$ joins two
distinct vertices when they share a denominator (as an integer).  The
\emph{seed} is $\sigma=(t,-2pt,-2pt)$; its connected component is the
\emph{seed component} $C_\sigma(p)$.  The \emph{seed distance}
$\dist(p)\in\{0,1,\dots\}\cup\{\infty\}$ is the graph distance from $\sigma$
to the nearest positive vertex of $C_\sigma(p)$.
\end{definition}

\begin{conjecture}[seed-component conjecture]\label{conj:seed}
For every prime $p\equiv1\pmod4$ the seed component $C_\sigma(p)$
contains a positive vertex.
\end{conjecture}

A positive vertex is an ES solution.  So Conjecture~\ref{conj:seed} implies
ES for all primes $p\equiv1\pmod4$; the primes $p\equiv3\pmod4$ are
classical.  Moreover Conjecture~\ref{conj:seed} is an
\emph{algorithmic} statement: starting from an object available for every
$p$, walk through the finite graph by exact divisor moves until positivity
is reached.  It was the working hypothesis of a pointwise programme.  This
paper explains why that programme is blocked.

\subsection{The main message}
The obstruction is an instance of a principle that is informally familiar
from Elsholtz--Tao \cite[\S1, remark after Prop.~1.6]{ET}: an argument proving
ES at $p$ must fail when $p$ is replaced by an odd square, which rules out
covering congruences and polynomial identities.  We make this precise for
the signed graph.

\begin{principle}[formal blindness]
Call an argument \emph{formal} if, along a residue class
$q\equiv q_0\pmod M$ of the parameter $q=(p-1)/24$, every integer it
inspects is the value at $q$ of a fixed polynomial expression, and every
decision it takes depends only on the class of $q$ and on the \emph{shape} of
the factorisations of these values (which irreducible polynomial factors
occur, with which exponents).  Then, at a \emph{formally generic} prime
(where all polynomial values met are primes), the argument sees only
\emph{formal} solutions.  In a square class there is no formally positive
solution (Schinzel; Bright--Loughran).  Hence, under Hypothesis~H, every
formal argument fails for infinitely many primes.
\end{principle}

Theorem~\ref{thm:closed} makes this precise for explorations of $G_p$, and
Proposition~\ref{prop:schinzel} is its unconditional core.  We prove it in
three cases: polynomial identities (Schinzel, unconditional), explorations
of bounded radius (Theorem~\ref{thm:unbounded}), and the whole seed
component (Theorems~\ref{thm:astra} and~\ref{thm:F}).  The conclusion
(\S\ref{sec:discussion}) is that a pointwise proof must control the
\emph{actual} factorisations of specific shifted integers, for instance by
using the size of~$p$.

\subsection{Results}
Status labels follow the campaign convention: \status{proved},
\status{conditional} (proved from a named hypothesis), \status{certified}
(an exact finite computation re-checked independently) and \status{evidence}
(finite data).  ``Verified numerically'' never means ``proved''.

\begin{introthm}[\status{proved}; \S\S\ref{sec:graph}--\ref{sec:short}]\label{thmA}
\leavevmode
\begin{enumerate}
\item A vertex is positive if and only if the two denominators of equal
$p$-adic valuation, with $p$ removed, have product a quadratic non-residue
modulo $p$ \emph{(Bright--Loughran \cite{BL}; Yamamoto for the positive
direction)}.
\item Every $p$-free denominator outside $[1,2t]$ lies in at most one vertex.
Every Type~II $p$-divisible denominator lies in at most two vertices, all of
one sign class.
\item $\dist(p)\le3$ if and only if one of three explicit divisor conditions
holds (Theorem~\ref{thm:depth3}).  Every prime $p\equiv1\pmod4$ outside
Mordell's six classes has $\dist(p)=2$.
\item $\#\{p\le N:\dist(p)>2\}\ll N(\log N)^{-11/2}$ and
$\#\{p\le N:\dist(p)>5\}\ll N(\log N)^{-10}$.
\end{enumerate}
\end{introthm}

\begin{introthm}[\status{conditional} on Hypothesis~H; Theorem~\ref{thm:unbounded}]\label{thmB}
For every $k$ there are infinitely many primes $p=24q+1$, with $q$ prime,
such that every vertex within distance $k$ of the seed is nonpositive.  Under
the Bateman--Horn conjecture there are $\gg_k N(\log N)^{-C_k}$ of them up to
$N$.  Hence no bounded-depth search from the seed proves ES.
\end{introthm}

\begin{introthm}[\status{conditional}; \S\ref{sec:whole}]\label{thmC}
\leavevmode
\begin{enumerate}
\item \emph{(Companion computation \cite{AS}; the stronger version.)}  There
are $159$ explicit linear forms $F_i(n)$, among them $q$ and $p=24q+1$, with
the following property.  For all sufficiently large $n$ at which every
$F_i(n)$ is prime, the seed component of $p$ has no positive vertex.  On an
admissible subprogression, $p$ nevertheless has an explicit positive
solution, which therefore lies outside $C_\sigma(p)$.  Under Dickson's
conjecture for these $159$ forms there are infinitely many such primes.
\item \emph{(This paper; independent confirmation.)}  Under Hypothesis~H for
an explicit family of $6402$ polynomials of degree at most two, infinitely
many primes $p=24q+1$ have a seed component with no positive vertex.  That
component consists of the values of $7883$ explicit formal vertices.
\end{enumerate}
In particular Conjecture~\ref{conj:seed} is false under either hypothesis.
\end{introthm}

The first part of Theorem~\ref{thmC} is due to the companion project
\cite{AS}, and it is stronger than the second part in two respects.  It
needs only Dickson's conjecture, the linear case of Hypothesis~H, for $159$
forms; and it exhibits an ES solution at the same primes, so its primes
are counterexamples to \emph{reachability}, not candidates against ES.  The
second part was obtained independently, by a different certificate (a
closure over a polynomial frame rather than an affine packet), and it was
re-verified by an independent engine.

\subsection{What is not claimed}
ES is not proved or disproved here.  Conjecture~\ref{conj:seed} is
\emph{open} unconditionally; no sterile seed component is known at an actual
prime, and the families of Theorem~\ref{thmC} are far beyond computational
reach.  Numerically every prime $p\equiv1\pmod4$ below $10^{12}$ has
$\dist(p)\le3$ (\S\ref{sec:evidence}).  This is \status{evidence}, and by
Theorem~\ref{thmB} the pattern is conditionally false.

\subsection{Credits}
The character dichotomy (Theorem~\ref{thm:BL}) is due to Bright and
Loughran \cite[Thms.~1.2, 1.5]{BL}; its positive direction goes back to
Yamamoto \cite{Yamamoto} (see \cite[App.~A]{BL}).  Finiteness of signed
solutions is \cite[Lemma~3.10]{BL}.  The labels and the Type~I chart are
Elsholtz--Tao coordinates \cite[(2.1), (2.6), (2.7), (2.18), (2.21)]{ET}.
Rigidity of $p$-divisible denominators between types was observed for
positive solutions by Jiang \cite[Thm.~3.2]{Jiang}; that preprint has since
been withdrawn.  The polynomial-identity obstruction is Schinzel's
\cite{Schinzel}; see \cite[\S1]{ET} and \cite[Cor.~1.4]{BL}.  The
half-dimension idea in Theorem~\ref{thm:sieve} is Dahan's
\cite[Lemma~4.2, Thm.~4.3]{Dahan}.  We found no prior source for the
signed graph, the seed, Theorem~\ref{thm:depth3} or
Theorems~\ref{thmB}--\ref{thmC}; Theorem~\ref{thmB} is best read as the
H-conditional form of the Elsholtz--Tao remark, and its novelty is
unconfirmed.

%======================================================================
\section{The signed graph}\label{sec:graph}

\subsection{Fibres and finiteness}
Fix a denominator $z$ of some vertex.  Since $p$ is odd, $4z\ne p$, so
$4/p-1/z=r/s$ with $\gcd(r,s)=1$, $s>0$ and $r\ne0$.

\begin{lemma}[fibres; \status{proved}]\label{lem:fibre}
The vertices containing $z$ are exactly the triples
\[
 \Bigl(z,\ \frac{s+D}{r},\ \frac{s+s^2/D}{r}\Bigr),
\]
where $D$ runs over the signed divisors of $s^2$ with $D\ne-s$ and
$r\mid s+D$.  The divisors $D$ and $s^2/D$ give the same vertex.  Since
$\gcd(D,r)=1$, integrality of the second entry implies that of the third.
\end{lemma}

\begin{proof}
$1/y+1/w=r/s$ is equivalent to $(ry-s)(rw-s)=s^2$.  Put $D=ry-s$.  Then
$D(s+s^2/D)=s(D+s)$, and $\gcd(D,r)=1$ because $D\mid s^2$ and
$\gcd(r,s)=1$.  The value $D=-s$ gives $y=0$.
\end{proof}

Every vertex has a positive denominator at most $3p/4$, since otherwise its
positive terms would sum to less than $4/p$.  By Lemma~\ref{lem:fibre}, $G_p$
is therefore finite \cite[Lemma~3.10]{BL}.

\subsection{Valuations, types and labels}

\begin{lemma}[valuations; \status{proved}]\label{lem:val}
Every vertex has exactly one or exactly two denominators divisible by $p$,
and each of them has $p$-adic valuation exactly~$1$.
\end{lemma}

\begin{proof}
Clearing denominators gives $4xyz=p(xy+yz+zx)$, so $p$ divides some
denominator.  If $p$ divided all three, multiplying by $p$ would write $4$ as
a sum of three reciprocals of nonzero integers; but such a sum is at most~$3$.  If
exactly one denominator $p^km$ is $p$-divisible, then
$1/(p^km)=4/p-1/x-1/y$ has valuation $-1$, so $k=1$.  Suppose $x$ is
$p$-free and $p^im,p^jn$ with $i,j\ge1$, $p\nmid mn$.  If $i<j$ and
$j\ge2$, then $1/(p^jn)=4/p-1/x-1/(p^im)$ would have valuation
$\ge-i>-j$, which is impossible.  So $i=j$.  If $i=j\ge2$ then
$|4/p-1/x|\le2/p^2$.  This forces $x>0$ and $x\le p^2/(4p-2)<p/2$, whence
$0<|4x-p|\le2x/p<1$, a contradiction.
\end{proof}

A vertex is of \emph{Type~I} if exactly one denominator is
$p$-divisible and of \emph{Type~II} if two are.  For a
$p$-divisible denominator $pm$ put $\lambda(pm)=4m-1\bmod p$.

\begin{lemma}[labels; \status{proved}]\label{lem:labels}
At a Type~I vertex the $p$-divisible denominator has $\lambda=0$.  At a Type~II
vertex $(x,pm,pn)$,
\[
 (4m-1)(4n-1)\equiv1,\qquad n(4m-1)\equiv m\pmod p .
\]
In particular both labels are nonzero, and a $p$-divisible denominator occurs
only in vertices of one type.
\end{lemma}

\begin{proof}
Multiply the equation by $p$ and reduce modulo $p$.  For $(x,y,pm)$ this
gives $4\equiv1/m$.  For $(x,pm,pn)$ it gives $4\equiv1/m+1/n$, that is,
$4mn\equiv m+n$.
\end{proof}

On positive points these labels are Elsholtz--Tao parameters: for Type~I,
$4m-1=pe$ is \cite[(2.1)]{ET}, and for Type~II the identity above is
\cite[(2.21)]{ET} reduced modulo~$p$.

\begin{lemma}[small anchor; \status{proved}]\label{lem:anchor}
Every vertex has a $p$-free denominator $x$ with $1\le x\le2t$.  The
$p$-free denominator of a Type~II vertex is in $[1,2t]$.
\end{lemma}

\begin{proof}
For Type~II, $(x,pm,pn)$ gives $p/x=4-1/m-1/n\ge2$, so $0<x\le p/2$.  For
Type~I, $(x,y,pm)$, the label gives $m\equiv-t\pmod p$, hence $m\le-t$ or
$m\ge3t+1$.  So $R:=4/p-1/(pm)$ is positive, with $R>4/p$ if $m<0$ and
$R\ge3/(3t+1)$ if $m>0$.  The positive terms among $1/x,1/y$ sum to at
least $R$, so one of $x,y$ is positive and at most $2/R$.  This bound is
less than $p/2$, respectively at most $2t+2/3$.
\end{proof}

In the positive case this is the Monks--Velingker bound
\cite[Thm.~2.1(i)]{MV}.

\subsection{The character dichotomy}

\begin{theorem}[Bright--Loughran; \status{proved} in \cite{BL}]\label{thm:BL}
Let $v$ be a vertex.  Let $A,B$ be the two denominators of $v$ with equal
$p$-adic valuation, each divided by $p^{v_p}$.  Then
\[
 v\ \text{is positive}\iff\leg{AB}{p}=-1 .
\]
\end{theorem}

\begin{proof}[Derivation from \cite{BL}]
By Lemma~\ref{lem:val} the valuation pattern is $(1,0,0)$ or $(0,1,1)$,
so exactly one pair has equal valuation.  Label the denominators
$u_1,u_2,u_3$ so that $u_2,u_3$ is this pair.  Then $u_2/u_3=A/B$ is a
$p$-adic unit and $v_p(u_1u_3)=1$ in both patterns.  By \cite[(3.1),
Lemma~3.4]{BL} the Hilbert symbol $(-u_1/u_3,-u_2/u_3)_p$ equals
$\leg{-u_2/u_3}{p}=\leg{-AB}{p}=\leg{AB}{p}$, since $p\equiv1\pmod4$.  For
$n=p$, \cite[Thm.~1.2]{BL} says that this symbol is $-1$ at positive
solutions, and \cite[Thm.~1.5]{BL} says that it is $+1$ at integer solutions
which are not positive.  The symbol is invariant under permuting the $u_i$
\cite[\S1]{BL}.
\end{proof}

The direction ``positive $\Rightarrow$ non-residue'' is \cite[Cor.~1.3]{BL},
which unifies Yamamoto's conditions \cite{Yamamoto}.  An elementary proof
by quadratic reciprocity, in the style of \cite[Prop.~1.6]{ET}, is also
possible; we do not need it.  By \cite[Thm.~1.6]{BL}, the Brauer group of
the Erd\H{o}s--Straus surface modulo constants is $\Z/2$, generated by the
class behind Theorem~\ref{thm:BL}.  So this character is the \emph{only}
Brauer--Manin invariant.  It identifies positivity exactly, but it cannot
create it.

\begin{corollary}[Type~II homogeneity; \status{proved}]\label{cor:typeII}
A Type~II vertex $(x,pm,pn)$ is positive if and only if
$\leg{\lambda(pm)}{p}=-1$.  Hence all vertices containing a given Type~II
denominator $pm$ have the same sign class.
\end{corollary}

\begin{proof}
By Lemma~\ref{lem:labels}, $n\equiv m\lambda(pm)^{-1}$, so
$\leg{mn}{p}=\leg{\lambda(pm)}{p}$.  Apply Theorem~\ref{thm:BL}.
\end{proof}

\subsection{The Type I chart}

\begin{lemma}[Type I chart; \status{proved}]\label{lem:chart}
Let $pm$ be the $p$-divisible denominator of a Type~I vertex and $x$ one of
its $p$-free denominators.  Write $x=t+a$; then $m=ph-t$ for an integer $h$.
Put $H=4h-1$ and $e=(4a-1)h-a$.  Then:
\begin{enumerate}
\item $e\ne0$, $e\mid x^2$, and the third denominator is $xm/e$;
\item $Hx=m+e$ and $(4a-1)H-4e=1$, so $e\mid x^2\iff e\mid m^2$;
\item conversely, if $e\ne0$ and $e\mid(t+a)^2$, then
$(t+a,\,p(ph-t),\,(t+a)(ph-t)/e)$ is a Type~I vertex.  In particular $t+a$
is automatically $p$-free;
\item if $x\ge1$, the vertex is positive if and only if $a\ge1$ and $h\ge1$.
\end{enumerate}
\end{lemma}

\begin{proof}
Since $4m\equiv1$, $m\equiv-t\pmod p$.  Next, $4m-1=pH$, so
$4/p-1/(pm)=H/m$, and $1/x+1/y=H/m$ gives $y=mx/(Hx-m)$.  Expanding,
$Hx-m=(4h-1)(t+a)-(ph-t)=(4a-1)h-a=e$.  Since $\gcd(e,H)=1$, the
conditions $e\mid mx=(Hx-e)x$, $e\mid x^2$ and $e\mid(Hx-e)^2=m^2$ are all
equivalent.  For (3): if $p\mid t+a$, then two denominators would be
$p$-divisible, one of them with label $\lambda=pH\equiv0$, contrary to
Lemma~\ref{lem:labels}.  For (4): $a,h\ge1$ give $e\ge2ah>0$.  Conversely
a positive vertex has $pm>0$, so $h\ge1$, and $1/x<4/p$, so $x>p/4$, that
is, $a\ge1$.
\end{proof}

On positive points, $e=a^2d$ in the coordinates of \cite[Prop.~2.2,
(2.6)--(2.7)]{ET}.  The chart is symmetric in the retained denominators $x$
and $pm$: a move retaining $x$ keeps $a$, and a move retaining $pm$ keeps
$h$.

\subsection{The seed, the hubs and the bridge}

\begin{lemma}[\status{proved}]\label{lem:seed}
\leavevmode
\begin{enumerate}
\item The denominator $-2pt$ occurs only in $\sigma$.  Hence every neighbour
of $\sigma$ contains $t$, and no vertex containing $t$, nor any positive
denominator $z\le t$, is positive.
\item The vertices containing $t$ are
$V_1(h)=(t,\ p(ph-t),\ -t(ph-t)/h)$ for signed $h\mid t^2$, together with
the Type~II vertices $(t,\ p(\delta-t),\ p(t^2/\delta-t))$ for signed
$\delta\mid t^2$, $\delta\ne t$.  The value $\delta=-t$ gives $\sigma$.
The third entry of $V_1(h)$ is negative for both signs of $h$.
\item The path
\[
 \sigma,\quad(t,-t(p+1),-pt(p+1)),\quad(2t,2t+1,-pt(p+1)),\quad(2t,2t,-pt)
\]
lies in $G_p$.  The vertices containing $-pt$ are
$(-pt,\ t+d,\ t+t^2/d)$ for signed $d\mid t^2$, $d\ne-t$.
\end{enumerate}
\end{lemma}

\begin{proof}
(1) A vertex containing $-2pt$ is Type~II (its label is $-8t-1\equiv1$),
and its other $p$-divisible denominator $pn$ has $4n-1\equiv1$, so
$n\equiv-2t\pmod p$.  So $|m|,|n|\ge2t$, and $p/x=4-1/m-1/n$ gives
$x>0$ and $|4-p/x|\le1/t$.  If $x\ge t+1$ then
$4-p/x\ge3/(t+1)>1/t$.  If $x\le t$ then $p/x-4\ge1/t$, with equality only
for $x=t$, and then $1/m+1/n=-1/t$ forces $m=n=-2t$.  A positive vertex has
all denominators $>p/4$.

(2) At $z=t$, $4/p-1/t=-1/(pt)$.  For Type~I, the chart with $a=0$ gives
$e=-h$.  For Type~II, $1/m+1/n=-1/t$ is $(m+t)(n+t)=t^2$.  If $h>0$ then
$ph-t>0$, and if $h<0$ then $ph-t<0$; either way $-t(ph-t)/h<0$.

(3) The reciprocal sums are checked directly, using $2t+1=(p+1)/2$.
Consecutive vertices share $t$, $-pt(p+1)$ and $2t$.  Finally,
$4/p+1/(pt)=1/t$, and Lemma~\ref{lem:fibre} applies with $r=1$, $s=t$.
\end{proof}

So the seed component contains the whole fibres of $t$ and of $-pt$.  These
fibres have $3\tau(t^2)$ and $\tau(t^2)$ vertices.  They are the two terms
of the unique signed two-term identity $4/p=1/t-1/(pt)$.

\subsection{Rigid fibres}

\begin{lemma}[Type~II fibres; \status{proved}]\label{lem:typeII}
A Type~II denominator $pm$ lies in at most two vertices.  If $|m|\ge2t$,
then in every vertex $(x,pm,pn)$ containing it, $n$ lies in $[-2t,2t]$.
\end{lemma}

\begin{proof}
Write the vertex as $(x,pm,pn)$, and put $K=4m-1$ and $A=4-1/m\in[3,5]$.
Then $x=p/(A-1/n)$ and $nK\equiv m\pmod p$.  The class of $n$ has exactly one
representative $n_0\in[-2t,2t]$.  Every other representative has
$|n|\ge(p+1)/2$, and then, with $\eta=2/(p+1)$,
\[
 \Bigl|x-\frac{pm}{K}\Bigr|=\frac{p/|n|}{A(A-1/n)}\le
 \frac{p\eta}{3(3-\eta)}=\frac{2p}{9p+3}<\frac14 .
\]
So all such vertices have the same $x$, the integer nearest to $pm/K$.  But $x$
determines $n$, so there are at most two vertices.  If $|m|,|n|\ge2t$, the
argument of Lemma~\ref{lem:seed}(1) shows that the vertex is $\sigma$, where
$n=-2t$.
\end{proof}

\begin{theorem}[outer anchors; \status{proved}]\label{thm:outer}
Every $p$-free denominator $z$ with $z<0$ or $z>2t$ lies in at most one vertex.
\end{theorem}

\begin{proof}
By Lemma~\ref{lem:anchor} a vertex containing $z$ is of Type~I,
$(z,x,pm)$, with $x$ $p$-free.  Each of $x$ and $m$ determines the vertex.
As in Lemma~\ref{lem:anchor}, $m\le-t$ or $m\ge3t+1$.  Put $R=4/p-1/z$,
$\delta_1=1/(pt)$ and $\delta_2=1/(p(3t+1))$.  Then
$1/x=R-1/(pm)\in[R-\delta_2,R+\delta_1]$.

\emph{Case $z<0$.}  Here $R>4/p>\delta_2$ and $R+\delta_1>1/t$.  So $x$ lies
in an interval of length
$(\delta_1+\delta_2)/((R-\delta_2)(R+\delta_1))<
\frac{1}{t(3t+1)}\cdot\frac{p(3t+1)}{12t+3}\cdot t=\frac13$.

\emph{Case $z\ge p$.}  Here $R\ge3/p$.  The length is at most
$p^2/((3t+1)(9t+2))<1$.

\emph{Case $2t<z<p$.}  Use the chart of Lemma~\ref{lem:chart} anchored at
$z$.  Put $A=4z-p\in[4t+3,3p-4]$ and $f=z^2/e$, a signed divisor of $z^2$.
Then $x=Hz^2/e-z$, and $AH=4e+1$ gives
\begin{equation}\label{eq:vieta}
 x+z=Hf,\qquad A(x+z)=4z^2+f .
\end{equation}
The vertex is positive if and only if $H\ge3$.  Then $0<f\le(x+z)/3\le2t$;
otherwise $H\le-1$ and $f<0$.  By \eqref{eq:vieta}, $x$ determines
the vertex.  Suppose that two vertices share $z$, with $x_1<x_2$.  Then
$f_2-f_1=A(x_2-x_1)\ge A$.
\begin{itemize}
\item Both positive: $0<f_1,f_2\le2t<A$, which is impossible.
\item Mixed, positive $P$ and nonpositive $N$: then $f_P>0>f_N$, so
$x_P>x_N$.  Also $|f_N|\le x_N+z\le6t$, so $A(x_P-x_N)\le8t<2A$.  Hence
$x_P=x_N+1$ and $|f_N|=A-f_P\ge2t+3$.  If $H_N\le-5$ then
$|f_N|\le6t/5$, which is too small.  So $H_N=-1$, i.e.\ $m_N=-t$, and
$1/x_N+1/z=1/t$.  Write $x_N=t+d$ and $z=t+u$ with $du=t^2$ and $d<t<u$.
Then $d=\lambda\alpha^2$, $u=\lambda\beta^2$ and $t=\lambda\alpha\beta$ with
$\alpha<\beta$; put $\gamma=\beta-\alpha$.  A computation from
\eqref{eq:vieta} gives
\begin{gather*}
 f_P=f_N+A=\lambda\gamma(4\beta-\gamma)-1,\\
 x_P+z=w+1,\qquad z^2=\lambda\beta^2w,\qquad w:=\lambda(\alpha+\beta)^2 .
\end{gather*}
Positivity needs $f_P\mid z^2$ and $f_P\mid x_P+z=w+1$.  So
$\gcd(f_P,w)=1$, and then $f_P\mid\lambda\beta^2$.  As
$f_P\equiv-1\pmod\lambda$, in fact $f_P\mid\beta^2$, say $\beta^2=kf_P$
with $k\ge1$.  Then $\beta$ is a root of
$X^2-4k\lambda\gamma X+k(\lambda\gamma^2+1)$.  So the quarter-discriminant
$k(\lambda\gamma^2\mathcal M-1)$, with $\mathcal M=4k\lambda-1\ge3$, is a
square $s^2$, and $s^2\equiv-k\pmod{\mathcal M}$ with
$\gcd(k,\mathcal M)=1$.  But $(-1|\mathcal M)=-1$.  If $2\mid k$ then
$\mathcal M\equiv7\pmod8$, so $(2|\mathcal M)=1$.  For the odd part $k'$
of $k$, reciprocity and $\mathcal M\equiv-1\pmod{k'}$ give
$(k'|\mathcal M)=(-1|k')(-1)^{(k'-1)/2}=1$.  So the Jacobi symbol
$(-k|\mathcal M)=-1$, and $-k$ is not a square modulo $\mathcal M$; this
is a contradiction.
\item Both nonpositive: $|f_1|=|f_2|+A(x_2-x_1)\ge4t+4$, so as above
$H_1=-1$, $x_1=t+d$, $z=t+u$.  Then
$|f_2|\le x_1+z-A=2t+d-3u+1<1$, which is impossible.\qedhere
\end{itemize}
\end{proof}

The cases $z<0$ and $z\ge p$, and the bound of two vertices for
$2t<z<p$, were found first; the Vieta--Jacobi step closes the remaining
case.  Theorem~\ref{thm:outer} was brute-force checked for all
$p\equiv1\pmod4$, $p<3\cdot10^4$, and independently for $p<400$ on all
$z\in[-4000,-1]\cup[2t+1,4000]$; no violation was found.

\begin{corollary}[where sign changes happen; \status{proved}]\label{cor:crossing}
If a positive and a nonpositive vertex share a denominator $z$, then either
$z$ is $p$-free with $t<z\le2t$, or $z=p(ph-t)$ is a Type~I denominator
with $h\ge1$.
\end{corollary}

\begin{proof}
$z$ lies in a positive vertex, so $z>p/4>t$.  If $z$ is $p$-free, then
$z\le2t$ by Theorem~\ref{thm:outer}.  If $z$ is $p$-divisible, then it is
of Type~I by Corollary~\ref{cor:typeII}, and $h\ge1$ because $z>0$.
\end{proof}

%======================================================================
\section{Short escapes, forced exits and sieve bounds}\label{sec:short}

A positive $p$-free denominator $z$ is \emph{productive} if its fibre
contains a positive vertex.  For $t<z\le2t$ put $q_z=4z-p>0$.  By
Lemma~\ref{lem:fibre}, applied with $r=q_z$ and $s=pz$, the denominator $z$
is productive if and only if some positive $D\mid z^2$ satisfies one of
\begin{align}
 D&\equiv-z\pmod{q_z}&&\text{(Type II exit } (z,\,p(z+D)/q_z,\,p(z+z^2/D)/q_z)),\label{eq:exitII}\\
 D&\equiv-1/4\pmod{q_z}&&\text{(Type I exit } (z,\,p(z+pD)/q_z,\,(pz+z^2/D)/q_z)).\label{eq:exitI}
\end{align}
This is the classical anchor criterion.

\begin{theorem}[escapes of length $\le3$; \status{proved}]\label{thm:depth3}
Let $p=4t+1$ be prime.
\begin{enumerate}
\item $\dist(p)\ge2$, and $\dist(p)\le2$ if and only if
\begin{equation}\label{eq:nine}
 \exists\,h>0,\ h\mid t^2,\ \exists\,D>0,\ D\mid(ph-t)^2:\quad
 D\equiv-h\pmod{4h-1}.
\end{equation}
\item $\dist(p)\le3$ if and only if \eqref{eq:nine}, (A) or (B) holds:
\begin{enumerate}
\item[(A)] \emph{(negative-bucket transfer)} there are $c>0$ with $c\mid t^2$,
and $d>0$ with $d\mid(pc+t)^2$ and $d\equiv-c\pmod{4c+1}$, such that, with
$a=(d+c)/(4c+1)$, $x=t+a$ and $w=x(pc+t)/d$, one of $x,w$ is productive.
The path is
$\sigma\to(t,-p(pc+t),-t(pc+t)/c)\to(x,-p(pc+t),w)\to{}$positive;
\item[(B)] \emph{(Type~II swap)} for some signed $\delta\mid t^2$,
$\delta\ne\pm t$, and some $m\in\{\delta-t,t^2/\delta-t\}$, the other
vertex $(x',pm,pn')$ of the bucket of $pm$ (Lemma~\ref{lem:typeII}) exists
and its $p$-free denominator $x'$ is productive.
\end{enumerate}
\end{enumerate}
\end{theorem}

\begin{proof}
The first layer $L_1$ consists of the vertices containing $t$
(Lemma~\ref{lem:seed}(1)), and all of them are nonpositive.

\emph{Distance two.}  Let $P$ be a positive neighbour of $V\in L_1$.  It
shares some $z\ne t$ with $V$, by Lemma~\ref{lem:seed}(1).  If
$V=V_1(h)$, then $z$ is not the negative entry, so $z=p(ph-t)$ with $h>0$,
and $P$ is a positive vertex of chart column $h$.  By
Lemma~\ref{lem:chart} this means that some $a\ge1$ has $e=(4a-1)h-a$
dividing $m^2$, $m=ph-t$.  Here $e\equiv-h\pmod{4h-1}$, and conversely
every positive $D\equiv-h$ is of the form $e$ with
$a=(D+h)/(4h-1)\ge1$.  This is \eqref{eq:nine}.  If $V$ is of Type~II, then
it is nonpositive.  So both of its labels are residues by
Corollary~\ref{cor:typeII}, and both of its $p$-divisible buckets are
nonpositive.

\emph{Distance three.}  Let $\sigma,V_1,V_2,P$ be a shortest path with
$V_1\in L_1$, $V_2$ nonpositive and not containing $t$.  Let
$z_1\ne t$ be shared by $V_1$ and $V_2$.  By Lemma~\ref{lem:labels},
$V_2$ has the type of the bucket of $z_1$.
\begin{itemize}
\item $V_1=V_1(h)$ and $z_1$ its negative entry: then $V_2=V_1$ by
Theorem~\ref{thm:outer}, which is impossible.
\item $V_1=V_1(h)$, $h>0$, $z_1=p(ph-t)$: chart $V_2$ in column $h$ from
its $p$-free denominator $x_2\in[1,2t]$ (Lemma~\ref{lem:anchor}).
Nonpositivity forces $a\le0$, so $x_2\le t$ and
$e=a(4h-1)-h<0$, and the third entry is negative.  So $P$ shares
$p(ph-t)$ with $V_2$, and then $P$ is adjacent to $V_1$, contradicting
minimality.
\item $V_1=V_1(-c)$, $c>0$, $z_1=-p(pc+t)$: chart $V_2$ in column $h=-c$.
If $a\le0$, then $e=c(1-4a)-a>0$, the third entry is negative and
$x_2\le t$, so $V_2$ has no positive neighbour.  If $a\ge1$ then
$d:=-e=(4c+1)a-c>0$, and both $p$-free entries are positive.  By
Lemma~\ref{lem:chart}(2), $d\mid(t+a)^2$ if and only if $d\mid(pc+t)^2$.
This is (A).
\item $V_1$ of Type~II and $z_1=pm$: then $V_2$ is the other vertex of that
bucket, and it is nonpositive (Corollary~\ref{cor:typeII}).  The label of
its second $p$-divisible denominator is a residue, so that bucket is
nonpositive too.  Hence $P$ contains $x'$.  This is (B).
\end{itemize}
Conversely, every path listed in (A) and (B) is legal.
\end{proof}

The proof was checked against brute-force layered BFS: for all
$p\equiv1\pmod4$ with $p<3\cdot10^4$, for every fifth such prime up to
$3\cdot10^5$, and independently for all $p<2500$.  The sets of positive
vertices reached by escapes of length exactly two and exactly three
coincide with those predicted by the theorem.

\begin{lemma}[exits need a fresh non-residue; \status{proved}]\label{lem:fresh}
\leavevmode
\begin{enumerate}
\item A positive exit of \eqref{eq:nine} with certificate $D$ requires
$\leg Dp=-1$.  Hence $ph-t$ has a prime factor $\ell$ with $\leg\ell p=-1$.
\item At an anchor $z>t$ with certificate $D\mid z^2$, the Type~I exit
\eqref{eq:exitI} requires $\leg Dp=-1$, and the Type~II exit
\eqref{eq:exitII} requires $\leg{zD}p=-1$.
\end{enumerate}
Every prime $\ell\mid t$ has $\leg\ell p=1$.  So the required non-residue
comes from a prime factor of a new shift ($ph-t$, $pc+t$, $t+a$, \dots),
not from the support of $t$.
\end{lemma}

\begin{proof}
(1) With $K=4h-1$ and $m=ph-t$, the $p$-free pair is
$y=(m+D)/K$, $w=(m+m^2/D)/K$, and $yw=m(m+D)^2/(DK^2)$.  Since
$m\equiv1/4\pmod p$, Theorem~\ref{thm:BL} gives the claim.  (2) The
$p$-free pair of \eqref{eq:exitI} has product $z^2(z+pD)/(Dq_z)$, and
$q_z\equiv4z\pmod p$.  For \eqref{eq:exitII} the quotient pair has product
$z(z+D)^2/(Dq_z^2)$.  Finally, if $\ell\mid t$ then $p\equiv1\pmod\ell$,
so $\leg\ell p=1$ by reciprocity (and $p\equiv1\pmod8$ if $2\mid t$).
\end{proof}

\begin{lemma}[forced exits; \status{proved}]\label{lem:forced}
\leavevmode
\begin{enumerate}
\item If $p\equiv5\pmod8$, then \eqref{eq:nine} holds with $h=1$, $D=2$.
\item If $p\equiv17\pmod{24}$, then \eqref{eq:nine} holds with $h=2$, $D=12$.
\item If $p\equiv1\pmod{24}$ and $p$ is a quadratic non-residue modulo $5$
or modulo $7$, then \eqref{eq:nine} holds as follows ($m=ph-t$):
\[
\begin{array}{llll}
\toprule
\text{class}&h&D&\text{reason}\\\midrule
p\equiv3\ (5)&1&5&t\equiv3\ (5),\ 5\mid3t+1,\ 5\equiv-1\ (3)\\
p\equiv2\ (5)&2&5&t\equiv4\ (5),\ 5\mid7t+2,\ 5\equiv-2\ (7)\\
p\equiv3\ (7)&6&63&m=23t+6,\ 21\mid m,\ 63\equiv-6\ (23)\\
p\equiv5\ (7)&3&63&m=11t+3,\ 21\mid m,\ 63\equiv-3\ (11)\\
p\equiv6\ (7)&4&m/14&m=15t+4,\ 14\mid m,\ m/14\equiv-4\ (15)\\\bottomrule
\end{array}
\]
\item \emph{(X-branch.)}  Let $d\mid t^2$, $d>0$, and $z=t+d$.  If some
positive $D\mid z^2$ has $D\equiv-1/4\pmod{4d-1}$, then $\dist(p)\le5$, via
the bridge of Lemma~\ref{lem:seed}(3), the vertex $(-pt,z,tz/d)$ and the
Type~I exit \eqref{eq:exitI} at $z$.
\end{enumerate}
\end{lemma}

\begin{proof}
(1) $t$ is odd, $m=3t+1$ is even, and $2\equiv-1\pmod3$.  (2) $t$ is even and
$t\equiv1\pmod3$, so $6\mid7t+2=2p-t$ and $12\mid m^2$, with
$12\equiv-2\pmod7$.  (3) In each row $6\mid t$, so $h\mid t^2$.  The stated
congruences are immediate (for the last row, $14\equiv-1\pmod{15}$).  (4)
$z$ is $p$-free by Lemma~\ref{lem:labels}, since $-pt$ has label $0$.  Its
modulus is $q_z=4d-1$.
\end{proof}

\begin{corollary}[Mordell's classes; \status{proved}]\label{cor:mordell}
Every prime $p\equiv1\pmod4$ outside the six classes
$p\equiv1,121,169,289,361,529\pmod{840}$ has $\dist(p)=2$.
\end{corollary}

These are the classical hard classes \cite{Mordell}, so the corollary adds
nothing to ES.  It shows that the seed reproduces Mordell's identities in
two moves, and that every question about the seed concerns only these
classes.

\begin{lemma}[half-dimension; adapted from {\cite[Lemma~4.2]{Dahan}}]\label{lem:half}
Let $K\ge3$, $c\in(\Z/K)^\times$, and let $N\ge1$ be coprime to $K$.  If no
positive divisor of $N^2$ is $\equiv c\pmod K$, then the set $R$ of residues
modulo $K$ of the prime factors of $N$ satisfies $R\cap cR^{-1}=\emptyset$.
In particular $|R|\le\varphi(K)/2$.
\end{lemma}

\begin{proof}
If $r,r'\in R$ with $rr'\equiv c$, choose primes $\ell\equiv r$ and
$\ell'\equiv r'$ dividing $N$.  If $\ell\ne\ell'$ then $\ell\ell'\mid N$;
otherwise $\ell^2\mid N^2$.  Either way some divisor of $N^2$ lies in the
class $c$.
\end{proof}

\begin{theorem}[\status{proved}, given a standard upper-bound sieve]\label{thm:sieve}
\begin{gather*}
 \#\{p\le N:\dist(p)>2\}\ll N(\log N)^{-11/2},\\
 \#\{p\le N:\dist(p)>5\}\ll N(\log N)^{-10} .
\end{gather*}
\end{theorem}

\begin{proof}
By Lemma~\ref{lem:forced}, $\dist(p)>2$ forces $p$ into one of the six
classes modulo $840$.  Then $6\mid t$, so every $h\mid36$ and every
$d\mid36$ divides $t^2$.  Consider the nine branches $N_h=ph-t$,
$K_h=4h-1$, $c_h=-h$ ($h\mid36$), and the nine branches $N_d=t+d$,
$K_d=4d-1$, $c_d=-1/4$ ($d\mid36$).  We have $\gcd(N_h,K_h)=1$, and
$\gcd(N_d,K_d)=1$ for $p>K_d$, because $4N_h=pK_h+1$ and $4N_d=p+K_d$.  If
$\dist(p)>2$, then each $h$-branch fails.  If $\dist(p)>5$, then each
$d$-branch also fails (Lemma~\ref{lem:forced}(4)).  By
Lemma~\ref{lem:half}, a failing branch $b$ confines the residues of the
prime factors of $N_b$ to a set $R_b$ with $|R_b|\le\varphi(K_b)/2$.

Fix one admissible family $(R_b)_b$; there are finitely many.  For a prime
$\ell$ outside a finite exceptional set, $\ell\mid N_h$ if and only if
$p\equiv-K_h^{-1}\pmod\ell$, and $\ell\mid N_d$ if and only if
$p\equiv-K_d\pmod\ell$.  For large $\ell$ these classes are pairwise
distinct and different from $0$, since $(1-4d)(4h-1)+1=4(h+d-4dh)\ne0$.
So the primes counted lie in a fixed class modulo $840$ and avoid, for
each $\ell\le N^{1/2}$, the class $0$ and every class of a branch $b$ with
$\ell\bmod K_b\notin R_b$.  By the prime number theorem in progressions,
this is a sieve problem of dimension
\[
 \kappa=1+\sum_b\Bigl(1-\frac{|R_b|}{\varphi(K_b)}\Bigr)\ge1+\frac B2 ,
\]
where $B=9$, respectively $18$, is the number of branches.  It has trivial
level of distribution, so an upper-bound sieve \cite{HR} gives
$\ll N(\log N)^{-\kappa}$.  Summing over the choices of $(R_b)$ proves the
theorem.
\end{proof}

The case $h=1$ alone gives the exponent $3/2$; compare
\cite[Thm.~4.14]{Dahan}.  These savings are only polylogarithmic.
Theorem~\ref{thm:unbounded} shows that, under Bateman--Horn, no bounded
depth can do better than polylogarithmic.  The data (\S\ref{sec:evidence})
suggest that the exponents are not sharp.

%======================================================================
\section{Formal genericity}\label{sec:formal}

\subsection{Hypotheses}
Let $\cP$ be the set of irreducible $g\in\Z[X]$ with content $1$ and positive
leading coefficient, and put $\Pp(X)=24X+1$.

\begin{hypothesis}[H; Schinzel--Sierpi\'nski \cite{SS}]
If $f_1,\dots,f_m\in\cP$ and no prime divides $f_1(y)\cdots f_m(y)$ for
all $y\in\Z$, then there are infinitely many $y$ with all $f_i(y)$ prime.
\end{hypothesis}

\emph{Dickson's conjecture} \cite{Dickson} is the case where all $f_i$ are
linear.  The \emph{Bateman--Horn conjecture} \cite{BH} predicts that
the number of such $y\le Y$ is asymptotic to
$C\,Y(\log Y)^{-m}/\prod_i\deg f_i$ for a constant $C>0$.  Each result below uses
only a finite, explicit instance.

\subsection{Formal frames}

\begin{definition}\label{def:frame}
A \emph{frame} $\cF=(\Lambda,E,q_0,\Se,\Sa)$ consists of a finite set
$\Lambda\supseteq\{2,3\}$ of primes, exponents $E_\ell\ge1$
($\ell\in\Lambda$), an integer $q_0$, and finite disjoint sets
$\Se,\Sa\subset\cP$ (\emph{entry} and \emph{auxiliary} polynomials) with
$X,\Pp\in\Se$.  For $g\in S=\Se\cup\Sa$ we require $g(q_0)\ne0$ and
$E_\ell>v_\ell(g(q_0))$, and we put
\[
 C_g=\prod_{\ell\in\Lambda}\ell^{v_\ell(g(q_0))},\qquad M=\prod_{\ell\in\Lambda}\ell^{E_\ell}.
\]
An integer $q$ is \emph{admissible} if $q\equiv q_0\pmod M$ and the numbers
$r_g=g(q)/C_g$, $g\in\Se$, are pairwise distinct primes outside $\Lambda$.
A \emph{formal integer} is $\varphi=c\prod_{g\in\Se}(g/C_g)^{e_g}$ with
$c\in\Z\setminus\{0\}$ and $e_g\ge0$.  Its value at $q$ is $\varphi(q)$,
and its \emph{sign} is the sign of $c$.  A \emph{formal vertex} is a triple of
formal integers with $1/a+1/b+1/c=4/\Pp$ in $\Q(X)$.  It is
\emph{formally positive} if all three constants are positive.
\end{definition}

\begin{lemma}[\status{proved}]\label{lem:values}
Let $q\equiv q_0\pmod M$ and $g\in S$.  Then $n_g(q):=g(q)/C_g$ is an
integer prime to every $\ell\in\Lambda$, and $n_g(q)\to\infty$ as
$q\to\infty$.  Suppose moreover that $C_X=C_\Pp=1$ and that $q$ is large
and admissible.  Then $q$ and $p=24q+1$ are prime.  The value of a formal
integer is a nonzero integer with the sign of $c$, and its $p$-adic valuation
is the exponent of $\Pp$.  Distinct formal integers have distinct values.
The values of a formal vertex form a vertex of $G_p$.
\end{lemma}

\begin{proof}
$v_\ell(g(q)-g(q_0))\ge E_\ell>v_\ell(g(q_0))$, so
$v_\ell(g(q))=v_\ell(C_g)$.  At admissible $q$ the value is
$c\prod r_g^{e_g}$, a factorisation into the distinct primes $r_g$.  These
exceed $|c|$ once $q$ is large, and $r_{\Pp}=p$.  Signs: every $g\in\cP$ has
positive leading coefficient.
\end{proof}

\begin{lemma}[H gives admissible values; \status{conditional}]\label{lem:H}
Assume that $C_X=C_\Pp=1$ and
\begin{itemize}
\item[(E3)] for every prime $\ell\notin\Lambda$ with
$\ell\le\sum_{g\in\Se}\deg g$, some residue class modulo $\ell$ is a root
of no $g\in\Se$.
\end{itemize}
Put $f_g(y)=g(My+q_0)/C_g$ for $g\in\Se$.  Then the $f_g$ lie in $\cP$ and
have no common fixed prime divisor.  Under Hypothesis~H for $\{f_g\}$ there
are infinitely many admissible $q$.  Under Bateman--Horn there are
$\gg N(\log N)^{-|\Se|}$ admissible $q$ with $24q+1\le N$.
\end{lemma}

\begin{proof}
The Taylor coefficients of $g(My+q_0)$ of positive degree are divisible by
$M$, and $C_g\mid M$, while $C_g\mid g(q_0)$; so $f_g\in\Z[y]$.  It is
irreducible with positive leading coefficient.  If $\ell\in\Lambda$, then
$f_g(y)\equiv f_g(0)\not\equiv0\pmod\ell$, since
$v_\ell(g(My+q_0)-g(q_0))\ge E_\ell>v_\ell(C_g)$.  This shows that $f_g$ is
primitive at $\ell$ and has no root modulo $\ell$.  If $\ell\notin\Lambda$,
then $y\mapsto My+q_0$ is invertible modulo $\ell$ and $\ell\nmid C_g$.  So
$f_g$ is primitive at $\ell$, and by (E3), or by degree counting when
$\ell>\sum\deg g$, $\prod f_g$ has a non-root modulo $\ell$.  For large
$y$ the prime values are pairwise distinct, because distinct elements of
$\cP$ are not proportional.  They lie outside $\Lambda$, because they are
$\ell$-units.
\end{proof}

\subsection{Formal fibres}
Let $Z$ be a formal integer.  Then $N:=4Z-\Pp\ne0$.  Cancel common
irreducible factors in $N/(\Pp Z)$ and write the result as $r/s$ with
\[
 r=c_r\prod_h(h/C_h)^{\beta_h},\qquad s=c_s\prod_{g\in\Se}(g/C_g)^{\alpha_g},
 \qquad\gcd(c_r,c_s)=1,\ c_s>0 .
\]
The polynomials of $s$ are those of $\Pp Z$, hence lie in $\Se$.  A
\emph{formal divisor} of $s^2$ is $D=\pm d\prod(g/C_g)^{j_g}$ with
$d\mid c_s^2$ and $0\le j_g\le2\alpha_g$.  For such $D$ consider
\begin{itemize}
\item[(A)] $h^{\beta_h}\mid D+s$ in $\Q[X]$ for every $h$ occurring in $r$;
\item[(B)] writing $D+s=k\prod_h h^{\beta_h}$ and
$K=k\prod_hC_h^{\beta_h}$, $v_\ell(K(q_0))\ge v_\ell(c_r)$ for every
$\ell\mid c_r$.
\end{itemize}
Let $\Fib(Z)$ be the set of $D\ne-s$ such that both $D$ and $s^2/D$ satisfy
(A) and (B).  For $D\in\Fib(Z)$ put $y_D=(D+s)/r$.

\begin{lemma}[formal fibres; \status{proved}]\label{lem:formalfibre}
Suppose that
\begin{itemize}
\item[(F1)] every irreducible factor of $N$ lies in $S$;
\item[(F2)] every prime of $c_r$ lies in $\Lambda$, and
$E_\ell\ge v_\ell(c_r)+v_\ell(\mathrm{den}(D+s))$ for all formal divisors
$D$ and all $\ell\mid c_r$.
\end{itemize}
Then for all large admissible $q$, the vertices of $G_p$ containing $Z(q)$
are exactly the triples $(Z(q),y_D(q),y_{s^2/D}(q))$ with $D\in\Fib(Z)$.
\end{lemma}

\begin{proof}
By Lemma~\ref{lem:fibre}, applied to the possibly non-reduced fraction
$r(q)/s(q)$ (the identity $(ry-s)(rw-s)=s^2$ does not need
reducedness), every vertex containing $z=Z(q)$ is
$(z,(s+D)/r,(s+s^2/D)/r)$ at $q$ for a divisor $D$ of $s(q)^2$.  By
Lemma~\ref{lem:values}, $s(q)=c_s\prod r_g^{\alpha_g}$ with distinct
primes $r_g>|c_s|$.  So its signed divisors are exactly the values of the
formal divisors, and $D=-s$ is excluded.  It remains to show that, for a
fixed formal divisor $D\ne-s$ and large admissible $q$, the number
$y_D(q)$ is an integer exactly when (A) and (B) hold.

\emph{(A) is necessary.}  Suppose $D+s=h^jk'$ with $j<\beta_h$ and
$h\nmid k'$.  Put $n=n_h(q)$, which is prime to $\Lambda$ by
Lemma~\ref{lem:values}.  If $y_D(q)\in\Z$, then $n^{\beta_h}$ divides the
integer $(D+s)(q)=y_D(q)\,r(q)=C_h^jn^jk'(q)$, because $r(q)=c_r\prod
n_{h'}(q)^{\beta_{h'}}$.  Let $d'$ be the denominator of $k'$.  By Gauss's
lemma it equals the denominator of $D+s$, so it is $\Lambda$-supported.
Then $d'k'\in\Z[X]$, and since $\pi\nmid d'C_h$ for every prime
$\pi\mid n$, we get $v_\pi(d'k'(q))\ge(\beta_h-j)v_\pi(n)\ge v_\pi(n)$.  Also
$v_\pi(h(q))\ge v_\pi(n)$.  Writing $uh+vd'k'=\mathrm{Res}(h,d'k')\ne0$ with
$u,v\in\Z[X]$, we get $n\mid\mathrm{Res}$, which fails once $n$ is large.

\emph{Given (A), (B) is exact.}  Here
$y_D(q)=K(q)\prod n_h^{\beta_h}/(c_r\prod n_h^{\beta_h})=K(q)/c_r$.  The
number $K(q)$ is an integer: its denominator is $\Lambda$-supported, and
$K(q)\prod n_h^{\beta_h}=(D+s)(q)\in\Z$ with $n_h$ prime to $\Lambda$.  So
$y_D(q)\in\Z$ if and only if $v_\ell(K(q))\ge v_\ell(c_r)$ for all
$\ell\mid c_r$.  By (F2), $K(q)\equiv K(q_0)$ to the precision needed.

The same applies to $s^2/D$.  Finally, $y_D\ne0$ for large $q$, because
$D+s\ne0$.
\end{proof}

A formal integer $Z$ is \emph{dead} if it is $p$-free (its $\Pp$-exponent
is $0$) and either its constant is negative or $Z-12X$ has positive leading
coefficient.  By Lemma~\ref{lem:values}, at large admissible $q$ its value is
a $p$-free denominator outside $[1,2t]$.  By Theorem~\ref{thm:outer} it then
lies in at most one vertex, so dead denominators never need to be expanded.

\subsection{Formal positivity is impossible in a square class}

\begin{proposition}[Schinzel; Bright--Loughran; \status{proved}]\label{prop:schinzel}
Let $\cF$ be a frame such that $24q_0+1$ is a nonzero square modulo every
$\ell\in\Lambda\setminus\{2,3\}$.  Assume that $\Lambda$ contains the primes
of $H_g:=24^{\deg g}g(-1/24)\in\Z\setminus\{0\}$ for every
$g\in\Se\setminus\{\Pp\}$.  Let $V$ be a formal vertex whose three
constants are $\Lambda$-supported up to sign.  Then $V$ is not formally
positive.
\end{proposition}

\begin{proof}
The class $24q_0+1\bmod 24M$ is a unit, so by Dirichlet's theorem there are
arbitrarily large $q\equiv q_0\pmod M$ with $p=24q+1$ prime.  By
Lemma~\ref{lem:values}, the values $n_g(q)$ are integers, so $V(q)$ is a
vertex of $G_p$ (no primality of $n_g(q)$ is needed here).  Since
$24q\equiv-1\pmod p$, we have $24^{\deg g}g(q)\equiv H_g\pmod p$.  So
$p\nmid n_g(q)$ for $g\ne\Pp$ once $p>|H_g|$, and the $p$-adic valuation of
each entry is its $\Pp$-exponent.  The equal-valuation pair, with $p$
removed, is therefore $\pm$ a product of primes of $\Lambda$ and of numbers
$n_g(q)$, $g\ne\Pp$.  For $\ell\in\Lambda$, $\ell\ge5$, we have
$\leg\ell p=\leg p\ell=\leg{24q_0+1}\ell=1$, while
$\leg2p=\leg3p=\leg{-1}p=1$ because $p\equiv1\pmod{24}$.  Also
$\leg{n_g(q)}p=\leg{g(q)}p=\leg{H_g}p=1$, because $24$, $C_g$ and $H_g$
are products of such primes.  By Theorem~\ref{thm:BL}, $V(q)$ is
nonpositive.  At large $q$ the signs of the entries are the signs of the
constants, so some constant is negative.
\end{proof}

A formal vertex is a polynomial solution of \eqref{eq:ES} in $n=24X+1$,
integral on the class $n\equiv24q_0+1\pmod{24M}$, and that class is a
square class.  Proposition~\ref{prop:schinzel} is thus a form of
Schinzel's theorem that no square class is solvable by polynomials with
positive leading coefficients \cite{Schinzel}, \cite[\S1]{ET}; the proof
is the standard reciprocity argument, here through Theorem~\ref{thm:BL}.

\subsection{Closed formal explorations}

\begin{theorem}[closed explorations are sterile; \status{conditional}]\label{thm:closed}
Let $\cF$ be a frame and $\cV$ a finite set of formal vertices such that
\begin{itemize}
\item[(E1)] $C_X=C_\Pp=1$ and $(6X,-12X\Pp,-12X\Pp)\in\cV$;
\item[(E2)] for every entry $Z$ of every member of $\cV$ that is not dead,
(F1) and (F2) hold, and for every $D\in\Fib(Z)$ the functions $y_D$,
$y_{s^2/D}$ are formal integers with $(Z,y_D,y_{s^2/D})\in\cV$;
\item[(E3)] as in Lemma~\ref{lem:H};
\item[(E4)] no member of $\cV$ is formally positive.
\end{itemize}
Then for every large admissible $q$ the seed component of $p=24q+1$ lies in
the value set of $\cV$ and has no positive vertex.  If $\cV$ is connected
under shared formal entries, then the seed component \emph{equals} this
value set.  Under Hypothesis~H for $\{f_g:g\in\Se\}$ there are
infinitely many such $p$, and under Bateman--Horn
$\gg N(\log N)^{-|\Se|}$ of them up to $N$.
\end{theorem}

\begin{proof}
The seed is the value of $(6X,-12X\Pp,-12X\Pp)$.  Let $v=V(q)$ with
$V\in\cV$, and let $v'$ be a neighbour sharing $z=Z(q)$, where $Z$ is an
entry of $V$.  If $Z$ is dead, its fibre is $\{v\}$.  Otherwise
Lemma~\ref{lem:formalfibre} and (E2) show that $v'$ is a value of a member
of $\cV$.  By induction on the distance, the seed component lies in the
value set.  By (E4) and Lemma~\ref{lem:values}, every value is nonpositive.
Formal vertices sharing a formal entry have values sharing a denominator,
which gives the equality.  The last sentence is Lemma~\ref{lem:H}.
\end{proof}

By Proposition~\ref{prop:schinzel}, (E4) holds automatically in a square
class once $\Lambda$ contains the primes of the constants and of the $H_g$.
In that case the only thing to check is the closure (E2).  This is the
precise form of the Principle for explorations of $G_p$: at an admissible
$q$, any exploration whose formal shadow is closed meets only formal
vertices, and none of them is positive.

\subsection{Bounded radius: unbounded seed distance}

\begin{lemma}[a generic base point; \status{proved}]\label{lem:generic}
There is $q^*=(q^*_\ell)\in\prod_\ell\Z_\ell$ such that (1) each
$24q^*_\ell+1$ is a nonzero square unit in $\Z_\ell$; (2) $g(q^*_\ell)\ne0$
for every $g\in\cP$ and every $\ell$; (3) for each $g\in\cP$,
$v_\ell(g(q^*_\ell))=0$ for all but finitely many $\ell$.
\end{lemma}

\begin{proof}
For $\ell=2,3$ every unit works, since $24u+1\equiv1$ modulo $8$,
respectively $3$.  Enumerate $\cP=\{g_1,g_2,\dots\}$.  For $\ell\ge5$ let
$n(\ell)$ be maximal with $\sum_{i\le n(\ell)}\deg g_i<(\ell-3)/2$.  There
are $(\ell-3)/2$ residues $r\not\equiv0$ with $24r+1$ a nonzero square.
Each primitive $g_i$ has at most $\deg g_i$ roots modulo $\ell$, so some such
$r$ is a root of none of $g_1,\dots,g_{n(\ell)}$.  Lift $r$ to a transcendental
element of $r+\ell\Z_\ell$.  Since $n(\ell)\to\infty$, (3) follows.
\end{proof}

For $g\in\cP$ put $C_g=\prod_\ell\ell^{v_\ell(g(q^*_\ell))}$.  Any frame
whose $\Lambda$ contains the primes of these $C_g$ ($g\in S$), and with
$q_0\equiv q^*_\ell\pmod{\ell^{E_\ell}}$, has the same constants $C_g$.

\begin{theorem}[unbounded seed distance; \status{conditional} on H]\label{thm:unbounded}
For every $k\ge0$ there is an explicit finite set $\Se^{(k)}\subset\cP$
with the following property.  Under Hypothesis~H for the associated
$\{f_g\}$, there are infinitely many primes $q$ with $p=24q+1$ prime such
that every vertex within distance $k$ of $\sigma$ is nonpositive.  In
particular $\dist(p)>k$.  Under Bateman--Horn there are
$\gg_kN(\log N)^{-|\Se^{(k)}|}$ such $p\le N$.
\end{theorem}

\begin{proof}
Fix $q^*$ as in Lemma~\ref{lem:generic}.  We build frames $\cF_j$
compatible with $q^*$ and sets $\cV_j$, by induction.  Start with
$\cV_0=\{(6X,-12X\Pp,-12X\Pp)\}$, $\Se=\{X,\Pp\}$, $\Sa=\emptyset$ and
$\Lambda=\{2,3\}$.  At step $j$, process every entry $Z$ of every member of
$\cV_j$, and enlarge the data as follows:
\begin{itemize}
\item add the factors of $4Z-\Pp$ to $\Sa$ (unless they are in $\Se$), and
add the primes of $c_r$ to $\Lambda$;
\item compute $\Fib(Z)$.  Condition (A) does not depend on $q$, and (B) is
decided by $q^*$.  Add the irreducible factors of every $y_D$ to $\Se$
(removing them from $\Sa$ if necessary), and the triples to $\cV_{j+1}$;
\item add to $\Lambda$ the primes of all new $C_g$, of all constants of
entries, of all $H_g$, and all primes $\le\sum_{\Se}\deg g$;
\item raise the $E_\ell$ to satisfy Definition~\ref{def:frame} and (F2),
and let $q_0$ be the CRT lift of $q^*$.
\end{itemize}
Each $y_D$ is a formal integer.  Its constant $c_y$ is rational.  By the
proof of Lemma~\ref{lem:formalfibre}, the value $y_D(q)=c_y\prod n_g(q)^{e_g}$
is an integer for every large $q\equiv q_0\pmod M$, and no primality is
needed for this.  A prime in the denominator of $c_y$ cannot lie in
$\Lambda$, since the $n_g(q)$ are prime to $\Lambda$.  It cannot lie outside
$\Lambda$ either: such a prime exceeds $\sum\deg g$, so it fails to
divide $\prod n_g(q)$ for some such $q$.  Admissibility
for $\cF_{j+1}$ implies admissibility for $\cF_j$.  By
Lemma~\ref{lem:formalfibre}, every neighbour of a value of $\cV_j$ is a
value of $\cV_{j+1}$.  So after $k$ steps every vertex within distance $k$
of $\sigma$ is a value of $\cV_k$, for large $\cF_k$-admissible $q$.  The
class $24q_0+1$ is a square class by Lemma~\ref{lem:generic}(1), so
Proposition~\ref{prop:schinzel} gives (E4) for $\cV_k$.  Condition (E3)
is vacuous, because the small primes lie in $\Lambda$.  Lemma~\ref{lem:H}
finishes the proof.
\end{proof}

\begin{remark}\label{rem:unbounded}
Only the entry polynomials need to take prime values; the auxiliary
factors of $4Z-\Pp$ need not.  The proof uses nothing about $\sigma$
except that its entries are formal.  So the same conclusion holds for any
bounded-radius exploration from formally given vertices, for instance any
``universal'' path built from divisors that exist for all $p$ in a residue
class.  Schinzel's theorem and \cite[Prop.~1.6]{ET} rule out arguments
that never factor.  Theorem~\ref{thm:unbounded} rules out, under H,
bounded-depth searches that do factor.  It is the H-conditional form of
the Elsholtz--Tao remark \cite[\S1]{ET}.  Related unconditional
obstructions for specific branches are \cite[Prop.~3.8,
Thm.~4.14]{Dahan}.  The family $\Se^{(k)}$ is large, so the theorem gives
no example within computational reach.
\end{remark}

\subsection{The whole seed component}\label{sec:whole}
Theorem~\ref{thm:unbounded} says nothing about the full component, which
is finite for each $p$ but whose radius is not bounded uniformly.  Two
independent constructions close it.  We state the stronger one first.

\begin{theorem}[companion computation \cite{AS}; \status{conditional} on Dickson]\label{thm:astra}
There are explicit integers $Q,M>0$ and $159$ explicit primitive linear forms
$F_i(n)=(\alpha_i(Q+Mn)+\beta_i)/C_i$ with the following property.  The forms
are $q=Q+Mn$, $p=24q+1$, and the primitive parts of $6q+a$ ($a\in A$,
$|A|=84$) and of $6(4h-1)q+h$ ($h\in H$, $|H|=75$), which correspond to the
shifted anchors $t+a$ and the Type~I quotients $ph-t$, with $q$ counted once.
\begin{enumerate}
\item The tuple $(F_i)$ is locally admissible.  There is a computable $N_0$
such that, if $n\ge N_0$ and all $F_i(n)$ are prime, then the seed component
of $p$ contains no positive vertex.
\item On the admissible subprogression $n\equiv507\pmod{857}$, where
$q\equiv713\pmod{857}$,
\[
 \Bigl(6q+7,\ (24q+1)(762q+32),\ \frac{(6q+7)(762q+32)}{857}\Bigr)
\]
is a positive solution.  It is the Type~I chart point $a=7$, $h=32$, $e=857$.
Under the hypotheses of (1) it lies outside the seed component.
\end{enumerate}
Consequently, under Dickson's conjecture for these $159$ forms, there are
infinitely many primes with a positive ES solution but a sterile seed
component.  An earlier version uses $800$ forms; it has the explicit
threshold $q>(Q+16)^{393600}$ and the positive point $a=10$, $h=33$,
$e=1277$ on $n\equiv1248\pmod{1277}$.
\end{theorem}

\begin{remark}[on Theorem~\ref{thm:astra}]
The proof in \cite{AS} combines a finite certificate with eventual
arguments.  The certificate is a packet of $302$ constant nodes and
$315{,}642$ signed divisor monomials, over a modulus of $287$ digits with
$91$ support primes.  The eventual arguments use polynomial division, with
one common threshold.  They also use ``one-prime'' fibre lemmas, which
exhaust the fibre of an affine anchor whose auxiliary partner is not
assumed prime.  This is why only the primary forms $t+a$ and $ph-t$ need
to be prime.  The construction also uses the outer-anchor bounds
(Theorem~\ref{thm:outer}) and Type~II rigidity.

\emph{Our verification status.}  We re-ran the companion's finite checker
on a copy of its data; it reports \texttt{VERIFIED\_PRIMARY\_PACKET\_EVENTUAL}
for the packet with SHA-256 \texttt{62a88612\dots c2f2}.  We checked the
positive solution in (2) symbolically.  We have \emph{not} independently
re-derived the eventual closure argument of \cite{AS}.  The proof of
Theorem~\ref{thm:astra} is the companion's.

Theorem~\ref{thm:astra} is stronger than Theorem~\ref{thm:F} below in three
ways.  It is linear: Dickson's conjecture, where one form is Dirichlet's
theorem, against quadratic prime values, which are not known even for a
single polynomial.  It is smaller: $159$ forms against $6402$, which also
improves the Bateman--Horn density.  And it shows that ES holds at the
primes concerned.
\end{remark}

\begin{theorem}[\status{conditional} on H; certificate \status{certified}]\label{thm:F}
There is an explicit frame $\cF$ with the following data: $|\Lambda|=2036$,
the largest prime in $\Lambda$ about $3.37\cdot10^{16}$, $E_\ell\le17$ and
$M\approx10^{12088}$.  The set $\Se$ consists of $6402$ polynomials ($5009$
linear, $1393$ quadratic, $\sum\deg=7795$), and $\Sa$ of $7119$ further
ones.  There is also a connected set $\cV$ of $7883$ formal vertices
satisfying (E1)--(E4).  Consequently, under Hypothesis~H for the $6402$
polynomials $f_g$, $g\in\Se$, infinitely many primes $p=24q+1$ ($q$ prime)
have a seed component that equals the value set of $\cV$ and contains no
positive vertex.  Under Bateman--Horn there are
$\gg N(\log N)^{-6402}$ such $p\le N$.
\end{theorem}

\begin{proof}
Apply Theorem~\ref{thm:closed} to the certificate described below.
\end{proof}

\emph{The certificate.}  The file \texttt{data/formal\_closure/certificate.json.gz}
contains $q_0\bmod\ell^{E_\ell}$ for every $\ell\in\Lambda$; the $13521$
polynomials of $S$, with their $C_g$ and entry/auxiliary flags; and the
$7883$ formal vertices, as triples $(c,\{(g,e_g)\})$.  The $7883$ vertices
have $9961$ non-dead denominators ($8412$ $p$-divisible and $1549$ $p$-free
anchors with $0<Z\le12X$) and $5726$ dead ones.  The certificate was found
by a fixed-point iteration.  A closure engine expands formal fibres from
the seed, and the class $q_0$ is adjusted at primes where the candidate
decisions (B) or the roots of $S$ would otherwise be non-generic.  It
stabilises after nine rounds.  The final closure uses no genericity
assumption: every prime of every constant $c_r$ is in $\Lambda$.

\emph{What is checked} (\status{certified}).  An exact verifier checks the
following:
\begin{itemize}
\item every $g\in S$ is primitive, irreducible and has positive leading
coefficient, and $C_g$ is correct;
\item every vertex satisfies the identity in $\Q[X]$, has formal-integer
entries, and has an entry with negative constant (E4);
\item the dead/anchor/$p$-divisible classification;
\item every non-dead fibre, recomputed and compared with $\cV$ (E2), and
the precision condition (F2);
\item condition (E3) for all $1271$ primes $\ell\notin\Lambda$ with
$\ell\le17631$;
\item connectivity of $\cV$.
\end{itemize}
The checks were run three times.  The verifier \texttt{formal2\_verify.py}
reported \texttt{OK} with $0$ mismatches.  It was re-run by the parent
session, together with \texttt{formal2\_verify\_extra.py}.  A hostile
reviewer then wrote a fibre engine from scratch that shares no code with the
generator.  It recomputed all $9961$ fibres: $533{,}011{,}471$ candidates
passed (A), and all fibres equal those of $\cV$, with no prime of a $c_r$
outside $\Lambda$ and no auxiliary factor outside $S$.  The reviewer also
brute-forced $16$ fibres at actual $q$ chosen to make the relevant values
prime, and all $16$ agree with the formal prediction.

\emph{Replay.}
{\small
\begin{verbatim}
F=data/formal_closure
uv run --with python-flint python scripts/formal2_verify.py \
    $F/lam_final.json $F/closure_final.json.gz
uv run --with python-flint python scripts/formal2_verify_extra.py \
    $F/lam_final.json $F/closure_final.json.gz
export PYTHONPATH=scripts
uv run --with python-flint python scripts/review_fc_global.py
uv run --with python-flint python scripts/review_fc_fibres.py \
    --all --jobs 4 --out /tmp/review_fc_fibres_all.json
\end{verbatim}}
The first command took $2$\,h\,$50$\,min on $14$ cores; the second takes
$4$ minutes and the third under one minute.  The from-scratch fibre
engine took $12$ minutes on $8$ cores.

\begin{remark}[characters in Theorem~\ref{thm:F}]\label{rem:Fsquare}
The class $q_0$ of Theorem~\ref{thm:F} is \emph{not} a square class: $24q_0+1$
is a non-residue modulo $870$ of the $2036$ primes of $\Lambda$.  These are
primes that were adjoined to fix the decisions (B) and the roots of $S$.
So (E4) is not an instance of Proposition~\ref{prop:schinzel}; it is
verified directly.  Theorem~\ref{thm:closed} itself never uses characters.
Only $6402$ of the $13521$ polynomials are required to be prime: the
auxiliary factors of $4Z-\Pp$ enter only through (A) and (B), which use
$n_h(q)\to\infty$ and $\Lambda$-freeness (Lemma~\ref{lem:formalfibre}).
\end{remark}

\begin{evidence}[actual $q$ in the certificate's class]\label{ev:realq}
Take $40$ primes $q\in[10^9,10^{11}]$ with $p=24q+1$ prime and
$q\equiv q_0$ modulo $2^4\cdot3^3\cdot5^2\cdot7\cdot11\cdot13$.  For each, an
exact layered BFS gives the following.
\begin{itemize}
\item The first layer ($80$ vertices) is always entirely formal.
\item Every component has a positive vertex at distance $2$.
\item On every escape path, the path leaves the formal set at a formal
denominator one of whose formal primes is not an actual prime.  In $31$
cases $g(q)/C_g$ is not even an integer; in $9$ it is composite.
\end{itemize}
So actual components follow the formal one exactly as far as the
admissibility conditions hold.  This illustrates the mechanism; it is not a
test of~H.
\end{evidence}

%======================================================================
\section{Further obstructions}\label{sec:further}

Each result below rules out a natural route to Conjecture~\ref{conj:seed}
or to a pointwise proof.

\subsection{Large denominators and Type~II buckets}
By Theorem~\ref{thm:outer}, Corollary~\ref{cor:typeII} and
Corollary~\ref{cor:crossing}, every edge of $G_p$ passes through a $p$-free
denominator in $[1,2t]$ or through a $p$-divisible one.  Type~II buckets
never mix signs and have at most two vertices.  Every sign change happens at
an anchor in $(t,2t]$ or at a Type~I bucket with $h\ge1$.
\emph{What this kills:} long hidden escape routes through large
denominators, and any argument that changes sign inside a Type~II bucket.
It is also what makes the formal closures of \S\ref{sec:formal} finite.

\subsection{Dead hubs and sterile components larger than the seed component}
Numerically, for typical $p$ the seed component is the largest component
(\S\ref{sec:evidence}).  This suggested a \emph{size conjecture}: every
component with more than $K(p)=O(\log p)$ vertices contains a positive
vertex.  Combined with ``the seed component is large'', it would have
implied Conjecture~\ref{conj:seed}.

\begin{lemma}[negative quadrant; \status{proved}]\label{lem:A}
A vertex containing a $p$-free denominator $x\in[1,t]$, or a Type~I
denominator $p(ph-t)$ with $h\le0$, is nonpositive.
\end{lemma}

\begin{proof}
$4/p-1/x<0$ in the first case, and $ph-t<0$ in the second.
\end{proof}

\begin{lemma}[dead-hub fibre; \status{proved}]\label{lem:B}
Let $k\ge2$, $\mathcal M=4k+1$, and suppose that $x=t-k\ge1$ has all prime
factors $\equiv1\pmod{\mathcal M}$.  Then the vertices containing $x$ are
exactly
\[
 V_d=\Bigl(x,\ -\frac{p(x-d)}{\mathcal M},\ \frac{px(x-d)}{\mathcal Md}\Bigr),
 \qquad d\mid x^2,\ 0<d<x .
\]
There are $(\tau(x^2)-1)/2$ of them; all are of Type~II and nonpositive.
\end{lemma}

\begin{proof}
Here $4x-p=-\mathcal M$.  By Lemma~\ref{lem:fibre} with $s=px$, the
vertices containing $x$ correspond to $f=\pm p^jd$ ($d\mid x^2$, $j\in\{0,1\}$)
with $f\equiv-x\pmod{\mathcal M}$; the $p$-divisible denominator is
$pm$ with $m=(x+f)/(-\mathcal M)$.  Modulo $\mathcal M$ we have
$d\equiv x\equiv1$ and $p=4x+\mathcal M\equiv4$.  So $f=d,pd,-pd$ would
need $\mathcal M\mid2,5,3$ respectively, and only $f=-d$ survives.  Then
$m=(d-x)/\mathcal M$, and $d\mapsto x^2/d$ swaps the two $p$-divisible
denominators.
\end{proof}

If, in addition, a guard prime keeps the descents of the spokes away from
small moduli, the component of $x$ is typically exactly the hub fibre plus
$2^r-1$ descent vertices.  It has size $(3^{r+1}-1)/2+2^r-1$ when $x$ is a
guard prime times $r$ primes.

\begin{evidence}[\status{certified} sterile components]\label{ev:sterile}
Each component below was exhausted and closed, and the result was
re-checked by an independent checker: every denominator's full fibre was
recomputed by two methods, and connectivity was verified.
\begin{center}
\begin{tabular}{rrrr}
\toprule
$p$ & sterile component & $\ln p$ & size$/\ln p$\\\midrule
$165426666927637$ & $398$ & $32.7$ & $12$\\
$1960717994383909$ (random $p$) & $339$ & $35.2$ & $9.6$\\
$475041241432047253$ & $3469$ & $40.7$ & $85$\\
$274159709010072908384347957$ & $30035$ & $61.9$ & $485$\\\bottomrule
\end{tabular}
\end{center}
At $p=274159709010072908384347957$ the seed component was also exhausted.
It has $10155$ vertices, $2811$ of them positive, while the sterile
component has $30035$.  In $13$ of $14$ comparable dead-hub primes the
sterile component was larger than the seed component.
\end{evidence}

\emph{What this kills:} every size threshold, and every
``giant component'' statement in the style of Bourgain--Gamburd--Sarnak with
the seed as the giant.  Variants of the same construction give certified
sterile components with $243$ distinct Type~I buckets, and with $146$
failing positive-capable denominators.  The asymptotic $O(\log p)$ form is
not formally refuted, since sterility of an infinite family is not proved.

\begin{lemma}[sign flip; \status{proved}]\label{lem:E}
Let $x=t+a$ with $a\ge1$ and $p\nmid x$, and suppose that $x$ has a prime
factor $\ell\equiv-1\pmod{4a-1}$.  If $x$ occurs in some vertex, then it
occurs in a positive one.  Likewise, if $h\ge1$ and $m=ph-t$ has a prime
factor $\ell\equiv-1\pmod{4h-1}$, then the Type~I bucket of $pm$ is empty or
contains a positive vertex.
\end{lemma}

\begin{proof}
Put $q=4a-1>0$.  The vertices containing $x$ are $(x,pm,pxm/f)$ with
$f=qm-x\mid px^2$ and $f\equiv-x\pmod q$, and they are positive exactly
when $f>0$.  Given such a vertex with $f=-p^jd<0$, replace $d$ by
$d'=\ell d$ if $v_\ell(d)<2v_\ell(x)$, and by $d/\ell$ otherwise.  Then
$d'\mid x^2$ and $d'\equiv-d\pmod q$, so $f'=p^jd'>0$ gives a positive
vertex; the third entry is integral because $\gcd(f',q)=1$.  The bucket
case is the same argument in the chart, using $e\mid m^2$ from
Lemma~\ref{lem:chart}(2).
\end{proof}

So in a sterile component every positive-capable anchor or bucket avoids
prime factors $\equiv-1$ modulo its own modulus.  This is a classical
divisor-class restriction, not a bound on size.  The certified dead hubs
contain few positive-capable denominators, all of large modulus.  The seed
component, by contrast, always contains the cheap tests $h\mid t^2$.  Any
replacement statement must weigh tests by their moduli, and this leads back
to the classical divisor-class problem.

\subsection{Parity and windmills}
Zagier's one-sentence proof \cite{Zagier} suggests looking for a finite set
with two involutions.

\begin{lemma}[\status{proved}]\label{lem:parity}
Let $\sigma_1,\sigma_2$ be involutions of a finite set $S$ with
$|\mathrm{Fix}\,\sigma_1|$ odd, and let
$\Phi:\mathrm{Fix}\,\sigma_2\to\{\text{positive vertices}\}$.  Then
$\sum_v|\Phi^{-1}(v)|$ is odd.  So any such argument exhibits an explicit
weight on positive solutions whose total is odd for every $p$.
\end{lemma}

The plain count is not such a weight: the number of positive solutions
is even for $p=97$ (it is $8$) and for $p=193$ (it is $6$).

\begin{proposition}[\status{proved}]\label{prop:windmill}
\leavevmode
\begin{enumerate}
\item \emph{(Repeated denominators.)} The only vertices with a repeated
denominator are $\sigma$ and $(2t,2t,-pt)$.  So coordinate permutations
have exactly two non-free orbits, and both are nonpositive.
\item \emph{(Fibre parity.)} Let $x>p/4$ be $p$-free, $r=4x-p$, $s=px$.  Let
$N^\pm$ be the number of positive $D\mid s^2$ with $D\equiv\mp s\pmod r$.
Then $N^+$ is even and $N^-$ is odd.  The fibre has $N^+/2$ positive
vertices and $(N^--1)/2$ vertices with exactly one negative entry.  So the
odd count in a fibre always certifies \emph{mixed-sign} vertices.
\item \emph{(Affine symmetries.)} The integral affine maps of $\Z^4$
preserving the Type~II polynomial $k(4abc-p)-a-b$ up to a constant factor
form a group of order $4$.  The same holds for the Type~I polynomial
$k(4abc-1)-p(a+b)$.  On positive points only the swap $a\leftrightarrow b$
survives.
\item \emph{(The natural odd set.)} The set of ordered positive solutions of
$1/x+1/y+2/z=4/p$ is finite and odd, but its part with $z$ even has exactly
$6f(p)$ elements, where $f(p)$ is the number of unordered positive
solutions of \eqref{eq:ES} with $n=p$.
\end{enumerate}
\end{proposition}

\begin{proof}
(1) If $2/x+1/z=4/p$ with $p\nmid x$, then $z=px/(2(2x-p))$ and
$2x-p=\pm1$, so $x=2t$ and $z=-pt$.  If the repeated entry is $pm$, then
$x=pm/(2(2m-1))$ forces $m=-2t$ and $x=t$.  (2) Complementation
$D\mapsto s^2/D$ preserves both classes.  Its unique positive fixed point
$D=s$ lies in the class $+s$, not $-s$, since $r\nmid2s$.  Positive $D$ in
the class $-s$ give positive vertices; negative divisors give mixed ones,
or $y=0$ when $D=-s$.  (3) The quartic part $4abck$ forces a monomial
linear part with entries $\pm1$.  The absence of cubic terms forces the
translation to vanish.  The linear part $-pk-a-b$ forces $k\mapsto\pm k$
(as $p\ne\pm1$), $\{a,b\}$ to be stable, and $c\mapsto c$.  (4) The swap of
$x,y$ fixes exactly the three solutions of $1/x+1/z=2/p$, all with $z$ odd.
On $z$ even, $(x,y,z)\mapsto\{x,y,z/2\}$ is $6$-to-$1$ onto ES solutions by
(1).
\end{proof}

\emph{What this kills.}  Fibre-level parity (2) and coordinate
permutations (1) give odd counts only of nonpositive objects.  No integral
affine windmill preserves the equation identically (3); unlike Zagier's
quadratic form $x^2+4yz$, the four-parameter surface is multilinear and has
no Vieta involutions.  Elsholtz \cite{Elsholtz} notes the same fragility
of windmill partitions for general forms.  A parity proof would need a new
explicit odd weight (Lemma~\ref{lem:parity}).
\status{evidence}: among the 732 primes $p\equiv1\pmod{24}$ below
$6\cdot10^4$, about $7000$ candidate weights were tested, and none
generalised beyond chance.  Pieces that preserve only the finite point set
are not excluded by (3).

%======================================================================
\section{Numerical evidence}\label{sec:evidence}

Everything in this section is \status{evidence}.  Finite computations do
not prove Conjecture~\ref{conj:seed}, and by Theorems~\ref{thm:unbounded}
and~\ref{thm:F} the patterns are conditionally misleading.

\begin{evidence}[seed distance below $10^{12}$]\label{ev:survey}
Every prime $p\equiv1\pmod4$ with $p<10^{12}$ has $\dist(p)\le3$; exactly
$44197$ of them have $\dist(p)=3$, and no prime with $\dist(p)\ge4$ is
known.  By Corollary~\ref{cor:mordell} only Mordell's six classes need
testing.  Their primes below $10^{12}$ were tested exhaustively by the
exact criterion \eqref{eq:nine}, and every survivor was re-checked through
the families (A), (B) of Theorem~\ref{thm:depth3}.
\begin{center}
\begin{tabular}{lrrr}
\toprule
range & primes tested & $\dist=3$ & $\dist\ge4$\\\midrule
$p<10^8$ & $179468$ & $70$ & $0$\\
$10^8\le p<10^{10}$ & $14036239$ & $1306$ & $0$\\
$10^{10}\le p<10^{11}$ & $114455512$ & $6196$ & $0$\\
$10^{11}\le p<10^{12}$ & $1046544177$ & $36625$ & $0$\\\bottomrule
\end{tabular}
\end{center}
Independently, a complete BFS covered all $p\equiv1\pmod4$ below
$3\cdot10^5$; only $297049$ has $\dist=3$.  Among the sparse inputs
$p=24q+1$ with $q$ prime and $q\le10^{13}$, all $1113907$ that fail
\eqref{eq:nine} escape at depth three, every one of them through family
(A).
\end{evidence}

\begin{evidence}[why nothing deeper is seen]\label{ev:anchors}
Among the survivors of \eqref{eq:nine} with $10^8<q\le10^{10}$, the
number of distinct productive (A)/(B) anchors has mean $14.7$ and minimum
$2$.  The mean grows with $p$ ($17.3$ near $q\approx5\cdot10^{12}$),
because the shifts $pc+t$ acquire more divisors.  So within reach
``distance $\le3$'' will keep looking true.  Theorem~\ref{thm:unbounded}
produces deep primes only from simultaneous primality of the family
$\Se^{(3)}$, whose Bateman--Horn density is far beyond reach.  The two
pictures are consistent.
\end{evidence}

\begin{evidence}[component sizes]\label{ev:sizes}
Complete enumeration was carried out for all $p\equiv1\pmod4$ below
$2\cdot10^6$, with samples up to $2^{31}$.  For typical $p$ the seed
component is the unique giant; its share of all vertices declines from
$0.58$ to $0.36$.  The mean per-prime maximal sterile size fits
$(\log_2p)^2/50$.  Nearly every large sterile component is a hub around a
negative-quadrant denominator (Lemma~\ref{lem:A}), which is exactly the
mechanism amplified in Lemma~\ref{lem:B}.
\end{evidence}

%======================================================================
\section{What a pointwise proof must do}\label{sec:discussion}

The results above locate the difficulty precisely.

\begin{enumerate}
\item \emph{Positivity is a character event.}  By Theorem~\ref{thm:BL} a
vertex is positive exactly when its equal-valuation pair has opposite
quadratic characters modulo $p$.  This is the only Brauer--Manin invariant
\cite[Thm.~1.6]{BL}, and it gives no obstruction to existence
\cite[Thm.~1.1]{BL}.  It cannot create a sign change; it only detects one.
\item \emph{Sign changes need fresh non-residues.}  By
Lemma~\ref{lem:fresh}, an exit requires a prime factor $\ell$ of a new
shift with $\leg\ell p=-1$.  Everything built from the support of $t$ is a
residue (a ``principal-genus trap'').
\item \emph{Formal arguments cannot supply them.}  If every number met
factors formally in a square class, every prime met is a residue
(Proposition~\ref{prop:schinzel}).  Under H this happens for infinitely
many primes, for any finite exploration (Theorems~\ref{thm:closed}
and~\ref{thm:unbounded}) and even for the whole seed component
(Theorems~\ref{thm:astra} and~\ref{thm:F}).  This is the Elsholtz--Tao
odd-square principle in graph form.
\item \emph{Global invariants do not help.}  Component size (\S5.2), parity
(\S5.3) and quadratic or quartic colourings give no invariant that forces
positivity.  Sterile components can be larger than the seed component.
\end{enumerate}

A pointwise proof of ES must therefore use information that a formally
generic prime does not have.  It must control the \emph{actual}
factorisation of specific shifted integers $t+a$, $ph-t$, \dots, in the
required classes, for every $p$.  This could come from the size of $p$
relative to the numbers met, from analytic input about prime factors of
shifted primes, or from a structure invisible to formal polynomial
families.  For the seed graph this is exactly the problem that the
exceptional-set machinery controls only on average: Vaughan-type bounds
$E(N)\ll N\exp(-c(\log N)^{2/3})$ and their refinements.  Conditionally on
H, Conjecture~\ref{conj:seed} is false.  By Theorem~\ref{thm:astra},
reachability fails even at primes where ES holds.  So the signed graph is
not a route to ES, and further work should return to the exceptional-set
problem or find genuinely non-formal input.

\section*{Acknowledgements}
% TODO(parent): settle the exact citation/authorship form of the companion
% computation [AS] (erdos-straus-astra) and of this campaign.
Theorem~\ref{thm:astra} and its proof are due to the companion computation
\cite{AS}, which reached the conditional disproof independently, first, and
in the stronger linear form.
\emph{[TODO: final citation and authorship form of \cite{AS} to be settled.]}

\appendix
\section{Replay}\label{app:replay}
{\raggedright
All commands run from the repository root; heavy runs were done under
\texttt{ulimit -v} and \texttt{timeout}.
\begin{itemize}
\item Base lemmas, hubs and fibre oracles: \texttt{scripts/pointwise\_seed\_check.py},
\texttt{pointwise\_fibres\_check.py}, \texttt{pointwise\_incidence.py 297049}.
\item Theorem~\ref{thm:depth3} validation: \texttt{scripts/depth3\_check.py},
\texttt{depth3\_validate.py 13 30000}; Lemma~\ref{lem:forced}:
\texttt{depth5\_branches.py 13 3000000}.
\item Survey (Evidence~\ref{ev:survey}): \texttt{scripts/depth3\_allp.cpp},
\texttt{depth3\_sieve.cpp}, \texttt{depth3\_batch.py}.
\item Theorem~\ref{thm:F}: the four commands in \S\ref{sec:whole}; the
fixed-point construction is \texttt{scripts/formal2\_iter.py}; actual-$q$
behaviour (Evidence~\ref{ev:realq}): \texttt{formal2\_realq.py}.
\item Theorem~\ref{thm:astra} (companion, read-only):
\texttt{scripts/pointwise\_primary\_packet\_check.py} in \cite{AS}.
\item Sterile certificates (Evidence~\ref{ev:sterile}):
\texttt{PYTHONPATH=scripts uv run --with python-flint python
scripts/sterile\_certificate\_check.py data/sterile/certs/*.json.gz}.
\item Parity scans (\S5.3): \texttt{scripts/windmill\_*.py},
\texttt{windmill\_singleton.cpp}.
\end{itemize}
\par}

{\raggedright
\begin{thebibliography}{99}
\bibitem{AS} Companion computation \emph{erdos-straus-astra}:
\texttt{SIGNED\_SEED\_COUNTEREXAMPLE.md}, \texttt{PRIMARY\_SEED\_PACKET.md},
\texttt{AFFINE\_SEED\_CLOSURE.md}, \texttt{ONE\_PRIME\_FIBRES.md},
data \texttt{signed-seed-primary-packet.json}, revision \texttt{a036064}
(1 October 2026).  \emph{[TODO(parent): final citation form.]}
\bibitem{BH} P.~T.~Bateman and R.~A.~Horn, A heuristic asymptotic formula
concerning the distribution of prime numbers, Math.\ Comp.\ 16 (1962),
363--367.
\bibitem{BL} M.~Bright and D.~Loughran, Brauer--Manin obstruction for
Erd\H{o}s--Straus surfaces, Bull.\ Lond.\ Math.\ Soc.\ 52 (2020), 746--761;
arXiv:1908.02526v2.
\bibitem{Dahan} B.~Dahan, Sieve dimension and search depth for the
Erd\H{o}s--Straus conjecture, $n\equiv1\pmod{24}$, arXiv:2608.24035v1 (2026).
\bibitem{Dickson} L.~E.~Dickson, A new extension of Dirichlet's theorem on
prime numbers, Messenger of Math.\ 33 (1904), 155--161.
\bibitem{Elsholtz} C.~Elsholtz, A combinatorial approach to sums of two
squares and related problems, in: Additive Number Theory, Springer, 2010,
115--140.
\bibitem{ET} C.~Elsholtz and T.~Tao, Counting the number of solutions to the
Erd\H{o}s--Straus equation on unit fractions, J.\ Aust.\ Math.\ Soc.\ 94
(2013), 50--105; arXiv:1107.1010v6.
\bibitem{HR} H.~Halberstam and H.-E.~Richert, Sieve Methods, Academic
Press, 1974.
\bibitem{Jiang} X.~Jiang, Counting Type I and Type II solutions of the
Erd\H{o}s--Straus equation, arXiv:2609.09204v1 (2026; withdrawn in v2).
\bibitem{MD} S.~Mihnea and B.~C.~Dumitru, Further verification and
empirical evidence for the Erd\H{o}s--Straus conjecture, arXiv:2509.00128
(2025).
\bibitem{MV} M.~Monks and A.~Velingker, On the Erd\H{o}s--Straus conjecture:
properties of solutions to its underlying diophantine equation, manuscript,
2008.
\bibitem{Mordell} L.~J.~Mordell, Diophantine Equations, Academic Press,
1969, Chapter~30.
\bibitem{Schinzel} A.~Schinzel, On sums of three unit fractions with
polynomial denominators, Funct.\ Approx.\ Comment.\ Math.\ 28 (2000),
187--194.
\bibitem{SS} A.~Schinzel and W.~Sierpi\'nski, Sur certaines hypoth\`eses
concernant les nombres premiers, Acta Arith.\ 4 (1958), 185--208.
\bibitem{Yamamoto} K.~Yamamoto, On the Diophantine equation
$4/n=1/x+1/y+1/z$, Mem.\ Fac.\ Sci.\ Kyushu Univ.\ Ser.\ A 19 (1965), 37--47.
\bibitem{Zagier} D.~Zagier, A one-sentence proof that every prime
$p\equiv1\pmod4$ is a sum of two squares, Amer.\ Math.\ Monthly 97 (1990),
144.
\end{thebibliography}
\par}

\end{document}
